\section*{Chapter \ref{ch:b3:complexity-classes} --- Time and Space Complexity}

\begin{solution}{exo:b3:complexity-classes:1}
(a)~A factor $d$ with $1 < d < N$: one division, polynomial in the number of digits of $N$. (b)~The cycle, as an order of the vertices: check that each vertex appears once and that consecutive vertices, the last and the first included, are adjacent. (c)~The $k$ vertices: check that they are distinct and that each of the $k(k - 1)/2$ pairs is an edge. (d)~The completed grid: check that it extends the given entries and obeys every rule, each rule looking at a few cells.
\end{solution}

\begin{solution}{exo:b3:complexity-classes:2}
Decision version: given the distances and a number $B$, is there a tour of length at most $B$? A tour has $n$ legs below $2^{32}$, so its length lies in $[0, n \cdot 2^{32}]$. Binary search on $B$ over that interval finds the least $B$ with a yes answer, which is the optimum, in about $32 + \log_2 n$ calls; each call is polynomial, and so is the total.
\end{solution}

\begin{solution}{exo:b3:complexity-classes:3}
(a)~$O(b + \log n)$ bits: the current maximum and an index. (b)~The same: the previous number and an index. (c)~$O(nb)$ bits for the copy. (d)~$O(n \log n)$ bits: a queue of up to $n$ vertex numbers and $n$ bits for the visited set.
\end{solution}

\begin{solution}{exo:b3:complexity-classes:4}
21, 1 and 21. \code{find_assignment} asks once at the start, then once per variable whatever the answer; when every variable must be false, each of the 20 trials answers no and the second branch is taken without a call (\code{test_self_reduction_uses_n_plus_one_calls_and_exercise_4}).
\end{solution}

\begin{solution}{exo:b3:complexity-classes:5}
The buggy \omterm{def:b3:complexity-classes:np}{verifier} checks only that the \omterm{def:b3:complexity-classes:np}{certificate} has $n$ entries and that consecutive entries are adjacent: 0, 1, 0, 1 has four entries and uses the edge $0 - 1$ three times, but never visits 2 or 3. The fix requires the \omterm{def:b3:complexity-classes:np}{certificate} to be a permutation of $0, \ldots, n - 1$, \code{sorted(order) == list(range(n))}, as \code{hamiltonian_path} does.
\end{solution}

\begin{solution}{exo:b3:complexity-classes:6}
\code{exercises/ex6_clique.py}: reject unless \code{chosen} has $k$ distinct vertices of the graph; put both directions of every edge in a set, $O(m)$; then look up each of the $k(k - 1)/2$ pairs, $O(k^2)$. The test is \code{test_exercise_6_clique}.
\end{solution}

\begin{solution}{exo:b3:complexity-classes:7}
Facts used: adding edges never makes a non-colourable graph colourable; a vertex adjacent to two vertices of a triangle must take the third one's colour.

Add a triangle $t_0, t_1, t_2$ to $G$; the new graph is 3-colourable exactly when $G$ is, and in any colouring the triangle uses all three colours. For each vertex $v$ in turn, and for $c = 0, 1$: join $v$ to the two triangle vertices other than $t_c$, and ask the oracle; if it says yes, keep the edges and go to the next vertex, otherwise remove them. If neither $c = 0$ nor $c = 1$ works, join $v$ to $t_0$ and $t_1$ without asking: since the graph was colourable before, some colouring gives $v$ one of the three triangle colours, and it was not that of $t_0$ or $t_1$. The invariant ``the current graph is 3-colourable'' holds throughout, at most $2n + 1$ calls are made, and at the end each $v$ is forced to the colour of one triangle vertex: reading these gives a 3-colouring of $G$.
\end{solution}

\begin{solution}{exo:b3:complexity-classes:8}
Let $V$ be a \omterm{def:b3:complexity-classes:np}{verifier} for $A$ running in time $p(n)$; it reads at most $p(n)$ symbols of the \omterm{def:b3:complexity-classes:np}{certificate}. On $x$, enumerate the \omterm{def:b3:complexity-classes:np}{certificates} of length at most $p(n)$ one after the other, in one buffer of $p(n)$ cells, run $V(x, c)$ on each, reusing its space, and accept if one is accepted. The space is $O(p(n))$ for the buffer plus $O(p(n))$ for a run of $V$, which cannot use more cells than steps: polynomial. The time is exponential, which does not matter.
\end{solution}

\begin{solution}{exo:b3:complexity-classes:9}
A call with $k > 0$ makes at most $2n$ calls with $k - 1$, two per midpoint, so $C(k) \le 1 + 2n\,C(k - 1)$ with $C(0) = 1$, and $C(k) \le (2n + 1)^k$. For $n = 16$, $k = 4$: $33^4 = 1{,}185{,}921$, against 88,741 printed for the unreachable query. The reachable query stops at the first midpoint that works, since \code{any} stops at the first true; and on either query \code{and} skips the second half when the first is false. On the unreachable query no midpoint works, so every midpoint is tried at every level.
\end{solution}

\begin{solution}{exo:b3:complexity-classes:10}
Facts used: the input is on a read-only tape and does not count; $f(n) \ge \log n$.

A configuration of the machine on $x$ is its state, the position of the input head ($O(\log n)$ bits), the position of the work head ($O(\log f(n))$ bits) and the work tape ($O(f(n))$ symbols): $O(f(n))$ bits, so there are $N = 2^{O(f(n))}$ configurations. The configuration graph has an edge $c \to c'$ when one step leads from $c$ to $c'$; whether it does is decided by looking at the transition table and at the input symbol under $c$'s input head, in space $O(f(n))$, without storing the graph. Make the machine erase its tape and move its heads home before accepting, so that there is one accepting configuration $c_{\mathrm{acc}}$. The machine accepts $x$ if and only if $c_{\mathrm{acc}}$ is reachable from the start configuration, and Savitch's recursion with $k = \lceil \log_2 N \rceil = O(f(n))$ decides it with $k + 1$ frames of $O(f(n))$ bits: $O(f(n)^2)$. Computing $k$ requires $f$ to be computable in space $f$, or trying $k = 1, 2, \ldots$ in turn.
\end{solution}

\begin{solution}{exo:b3:complexity-classes:11}
Take $f(n) = 2^n$. The time hierarchy theorem gives a problem $A$ decided in time $O(2^n \cdot n)$, so $A \in \mathrm{TIME}(2^{n^2}) \subseteq \mathrm{EXP}$, but in no time $o(2^n)$. Every problem in P is in $\mathrm{TIME}(n^c)$ for some $c$, and $n^c = o(2^n)$. So $A$ is not in P: $\mathrm{P} \ne \mathrm{EXP}$.
\end{solution}

\begin{solution}{exo:b3:complexity-classes:12}
\code{exercises/ex12_bipartite.py}. A walk from $v$ back to $v$ of odd length exists if and only if the graph has an odd cycle, since an odd closed walk contains an odd cycle; so the graph is 2-colourable if and only if no $(v, 0)$ reaches $(v, 1)$ in the parity graph. That graph has $2n$ vertices, and an edge $(x, p) \to (y, 1 - p)$ for each edge $xy$; the function \code{adjacent} reads it from the edge set without building it. The recursion is $\lceil \log_2 2n \rceil + 1$ frames deep, each holding six small numbers, $O(\log n)$ bits: $O(\log^2 n)$ space, the input aside (\code{test_exercise_12_bipartite}).
\end{solution}

\begin{solution}{pb:b3:complexity-classes:1}
\textbf{1.} $O(n + m)$: each vertex enters the queue once and each edge is looked at once. For the town, 30 places and 105 directed edges: about 135 steps.

\textbf{2.} A visited set of 30 places and a queue of at most 30: about 60 numbers of 5 bits, 300 bits; with the visited set as one bit per place, 30 bits plus the queue.

\textbf{3.} BFS must not visit a place twice, and in the worst case any place seen may come back; it keeps the whole frontier and the set of places seen, which grow with the map.

\textbf{4.} No: on a path the stack holds all $n$ vertices.

\textbf{5.} 4~GB for the queue in the worst case, plus 125~MB for a visited bitmap.

\textbf{6.} $k = \lceil \log_2 30 \rceil = 5$: a path between two places has at most 29 roads, fewer than $2^5 = 32$.

\textbf{7.} $u$, $v$, $w$ (5 bits each for 30 places) and $k$ (3 bits): 18 bits.

\textbf{8.} $5 + 1 = 6$, as printed.

\textbf{9.} $6 \times 18 = 108$ bits, against about 300 for BFS. Small maps hide the difference: $O(\log^2 n)$ against $O(n \log n)$.

\textbf{10.} $k = 30$, 31 frames of $30 + 30 + 30 + 5 = 95$ bits: about 2,950 bits, under 400 bytes, against gigabytes.

\textbf{11.} The second half reuses the memory of the first; evaluating both at once would double the memory at every level and make it linear in $n$ again.

\textbf{12.} $C(0) = 1$ and $C(k) \le 1 + 2n\,C(k - 1)$: each of $n$ midpoints may need two calls.

\textbf{13.} By induction: $1 + 2n(2n + 1)^{k - 1} \le (2n + 1)^{k - 1} + 2n(2n + 1)^{k - 1} = (2n + 1)^k$.

\textbf{14.} $\log_2 31 \approx 4.95$, $\log_2 820 \approx 9.68$, $\log_2 88{,}741 \approx 16.44$, against 4, 9 and 16: $\log_2 C \approx (\log_2 n)^2 + 0.5$, so $C \approx 1.4 \, n^{\log_2 n}$.

\textbf{15.} $n = 32$: $2^{25.5} \approx 4.5 \times 10^7$ calls, about 45 seconds. $n = 1{,}024$: $2^{100.5} \approx 1.8 \times 10^{30}$ calls, far longer than the age of the universe.

\textbf{16.} $n^{\log_2 n} = 2^{(\log_2 n)^2}$: the exponent $(\log_2 n)^2$ outgrows every $c \log_2 n$, so no polynomial bounds it, and stays below $n$, so it is less than $2^n$.

\textbf{17.} The place is reachable, and \code{any} stops at the first midpoint that works. Midpoint $u$ itself works at once for the first half, $R(u, u, k - 1)$, so the recursion soon reduces to shorter questions; the expensive case is a query that fails.

\textbf{18.} BFS on the town: it is fast and its memory is tiny anyway. A Savitch-style recursion only where memory is the constraint and the graph is given implicitly, by a function returning neighbours.

\textbf{19.} Nobody knows. Directed reachability is complete for NL under \omterm{def:b3:complexity-classes:logspace}{log-space reductions}, so such an algorithm would give $\mathrm{L} = \mathrm{NL}$.

\textbf{20.} A nondeterministic walk guesses the next vertex at each step and keeps only the current vertex and a counter up to $n$: $O(\log n)$ bits.

\textbf{21.} Space can be reused: Savitch's two halves share the same cells. Time cannot: a deterministic simulation of a guessing machine seems to need to try the guesses one after the other, and the only known bound is exponential, $\mathrm{NP} \subseteq \mathrm{EXP}$.

\textbf{22.} Its vertices are the configurations, its edges the single steps; the machine accepts if and only if a path leads from the start to the accepting configuration. A machine using space $s$ gives a maze of $2^{O(s)}$ rooms.

\textbf{23.} A nondeterministic machine using space $n^c$ is simulated deterministically in space $n^{2c}$, still polynomial: $\mathrm{NPSPACE} \subseteq \mathrm{PSPACE}$, and the other inclusion is immediate.

\textbf{24.} BFS spends memory proportional to the graph and saves time, $O(n + m)$; Savitch's recursion spends time, $n^{O(\log n)}$, and saves memory, $O(\log^2 n)$ bits. On the town both answer at once; the theorem matters when the graph is a configuration graph, far too large to store.
\end{solution}
